\subsection{רקע כללי בלמידה עמוקה}

למידה עמוקה (\textenglish{Deep Learning}) היא תת־תחום של למידת מכונה
(\textenglish{Machine Learning}), המבוסס על שימוש ברשתות נוירונים
מלאכותיות בעלות שכבות רבות. מטרת המערכת היא ללמוד ייצוגים מורכבים של
נתונים מתוך דוגמאות, ללא צורך בהגדרה ידנית של המאפיינים על ידי המתכנת.

\subsubsection{רשתות נוירונים ותהליך ה־\textenglish{Forward Pass}}

רשת נוירונים היא מודל חישובי השואב השראה מפעולת המוח האנושי. היא בנויה
מצמתים (נוירונים) המאורגנים בשכבות:
\begin{itemize}
\item \textbf{שכבת קלט (\textenglish{Input Layer}):} מקבלת את הנתונים
הגולמיים.
\item \textbf{שכבות חבויות (\textenglish{Hidden Layers}):} שכבות אלו
מבצעות חישובים מתמטיים הכוללים כפל של הקלט ב"משקולות"
(\textenglish{Weights}), הוספת הסטה (\textenglish{Bias}) והפעלת
פונקציית אקטיבציה.
\item \textbf{שכבת פלט (\textenglish{Output Layer}):} מחזירה את תוצאת
המודל.
\end{itemize} תהליך ה־\textbf{\textenglish{Forward Pass}} הוא השלב שבו
המידע זורם מהקלט דרך השכבות השונות עד לקבלת החיזוי הסופי.

\subsubsection{פונקציות אקטיבציה}

פונקציות אקטיבציה מכניסות אי־ליניאריות למודל. ללא אי־ליניאריות, גם רשת
בעלת שכבות רבות הייתה שקולה למודל ליניארי פשוט, ולכן לא הייתה מסוגלת
ללמוד תבניות חזותיות מורכבות.

בפרויקט זה נעשה שימוש בפונקציות אקטיבציה ממשפחת \textenglish{ReLU}
(\textenglish{Rectified Linear Unit}), המוגדרת כך:
\[ f(x) = \max(0, x)
\] פונקציה זו יעילה במיוחד לרשתות עמוקות, משום שהיא פשוטה לחישוב
ומסייעת להפחתת בעיית דעיכת הגרדיאנט (\textenglish{Vanishing
Gradient}).

\subsubsection{תהליך הלמידה ואופטימיזציה}

כדי שהרשת תוכל להשתפר, עליה לעבור תהליך איטרטיבי של למידה המורכב
מהשלבים הבאים:
\begin{itemize}
\item \textbf{פונקציית הפסד (\textenglish{Loss Function}):} פונקציה
המחשבת את גודל הטעות של המודל בסוף ה־\textenglish{Forward Pass}. ככל
שערך ההפסד נמוך יותר, כך המודל מדויק יותר.
\item \textbf{מורד הגרדיאנט (\textenglish{Gradient Descent}):} זוהי
האסטרטגיה המרכזית להקטנת השגיאה. אנו מתייחסים לפונקציית ההפסד כאל
"נוף" הררי ומחפשים את הנקודה הנמוכה ביותר בו. הגרדיאנט (הנגזרת) מסמן
לנו את כיוון העלייה, ולכן אנו צועדים בכיוון ההפוך לו כדי "לרדת" במדרון
לעבר מינימום השגיאה.
\item \textbf{פעפוע לאחור (\textenglish{Backpropagation}):} האלגוריתם
המחשב את הגרדיאנטים בפועל. באמצעות "כלל השרשרת" (\textenglish{Chain
Rule}), הרשת מחשבת את הנגזרת של פונקציית ההפסד ביחס לכל משקולת ברשת,
מהסוף (הפלט) חזרה להתחלה (הקלט).
\item \textbf{אופטימייזר (\textenglish{Optimizer}):} האופטימייזר הוא
הרכיב המבצע את עדכון המשקולות בפועל. בפרויקט זה נעשה שימוש באלגוריתם
\textbf{\textenglish{AdamW}}, המשלב קצב למידה אדפטיבי (התאמת גודל הצעד
לכל משקולת), מומנטום (זיכרון של כיווני העדכון הקודמים) ודעיכת משקולות, מה שמאפשר
למידה מהירה ויציבה יותר.
\end{itemize}

\subsubsection{רגולריזציה, \textenglish{Overfitting}
ו־\textenglish{Underfitting}}

המטרה של אימון רשת איננה רק להגיע לביצועים טובים על דוגמאות האימון,
אלא להשיג \textbf{הכללה} (\textenglish{Generalization}) לנתונים חדשים.

שני מצבים בעייתיים במיוחד הם:
\begin{itemize}
\item \textbf{\textenglish{Overfitting} (התאמת יתר):} המודל לומד היטב
את נתוני האימון, אך מתקשה להכליל לנתונים שלא ראה.
\item \textbf{\textenglish{Underfitting} (התאמת חסר):} המודל פשוט מדי,
או לא אומן מספיק, ולכן אינו מצליח ללמוד גם את הדפוסים הבסיסיים.
\end{itemize}

כדי להתמודד עם \textenglish{Overfitting} משתמשים לעיתים בכמה טכניקות
רגולריזציה:
\begin{itemize}
\item \textbf{\textenglish{Dropout}:} השבתה אקראית של נוירונים בזמן
האימון כדי למנוע הסתמכות יתר על מאפיינים מסוימים.
\item \textbf{\textenglish{Batch Normalization}:} נרמול האקטיבציות
בתוך הרשת, המשפר את יציבות האימון ולעיתים גם את ההכללה.
\item \textbf{\textenglish{Early Stopping}:} עצירת האימון כאשר ביצועי
הוולידציה מפסיקים להשתפר.
\end{itemize}

בפרויקטים של ראייה ממוחשבת גם אוגמנטציה (\textenglish{Data
  Augmentation}) ופיצול נכון של הנתונים למערכי אימון, וולידציה ובדיקה
מסייעים להפחתת התאמת יתר.

\subsection{זיהוי אובייקטים (\textenglish{Object Detection})}

זיהוי אובייקטים הוא משימה בראייה ממוחשבת שבה המודל צריך לענות על שתי
שאלות בו־זמנית: \textbf{מה} מופיע בתמונה, ו\textbf{איפה} הוא מופיע.
בפרויקט שלי משימה זו מופיעה פעמיים: פעם אחת לזיהוי הרובוטים על המגרש,
ופעם נוספת לזיהוי הספרות שעל הבאמפר של כל רובוט.

\subsubsection{ההבדל בין \textenglish{Classification}
ל־\textenglish{Detection}}

במשימת סיווג (\textenglish{Classification}) המודל מקבל תמונה ומחזיר
מחלקה. לעומת זאת, בזיהוי אובייקטים (\textenglish{Detection}) המודל
צריך גם לאתר כל אובייקט בתוך התמונה באמצעות תיבה תוחמת
(\textenglish{Bounding Box}).

בפרויקט שלי ההבדל הזה קריטי, משום שבכל פריים יכולים להופיע כמה רובוטים
במקביל. לכן לא מספיק לדעת שיש בתמונה "רובוט", אלא צריך לדעת היכן נמצא
כל רובוט כדי לעקוב אחריו לאורך זמן ולשייך אותו לקבוצה הנכונה.

\subsubsection{רשתות קונבולוציה (\textenglish{CNN})}

הבסיס למודלי זיהוי מודרניים הוא רשתות קונבולוציה
(\textenglish{Convolutional Neural Networks}). רשתות אלו משתמשות
במסננים (\textenglish{Filters}) הסורקים את התמונה ולומדים לזהות תחילה
מאפיינים פשוטים יחסית, כמו קצוות וקווים, ובהמשך מאפיינים מורכבים יותר,
כמו חלקי אובייקט או צורות שלמות.

פלט הביניים של שכבות הקונבולוציה נקרא \textenglish{Feature
Maps}. לעיתים משלבים גם שכבות \textenglish{Pooling}, שמקטינות את הממד
המרחבי של המפות, מפחיתות עומס חישובי, ומבליטות מאפיינים חשובים.

\subsubsection{ארכיטקטורת \textenglish{YOLO}}

בפרויקט זה נעשה שימוש במודלי \textenglish{YOLO} (\textenglish{You Only
Look Once}), השייכים למשפחת הגלאים החד־שלביים (\textenglish{One-Stage
Detectors}). בניגוד למשפחות ישנות יותר כמו \textenglish{R-CNN}
ו־\textenglish{Faster R-CNN}, שמפרידות בין שלב הצעת אזורים לשלב
הסיווג, \textenglish{YOLO} מבצע את כל החיזוי במעבר יחיד על התמונה. לכן
הוא מתאים במיוחד ליישומים שבהם המהירות חשובה.

נהוג לחלק את הארכיטקטורה לשלושה חלקים:
\begin{itemize}
\item \textbf{\textenglish{Backbone}:} החלק שמחלץ מאפיינים חזותיים מן
התמונה.
\item \textbf{\textenglish{Neck}:} החלק שמחבר מאפיינים מכמה קני מידה,
כדי לאפשר זיהוי של אובייקטים בגדלים שונים.
\item \textbf{\textenglish{Head}:} החלק שמחזיר את התיבות, ציוני
הביטחון והמחלקות החזויות.
\end{itemize}

בפרויקט שלי משתמשים בגישה זו הן לזיהוי רובוטים יחסית גדולים בתוך תמונת
המגרש, והן לזיהוי ספרות קטנות יותר מתוך חיתוכים של רובוט בודד.

\subsubsection{\textenglish{Anchor Boxes} ו־\textenglish{Non-Max
    Suppression}}

אחד הרעיונות המרכזיים במודלי זיהוי קלאסיים הוא \textenglish{Anchor
Boxes}: תיבות ייחוס בגדלים וביחסי ממדים שונים, שמהן המודל לומד לסטות
כדי להתאים את התיבה בפועל לאובייקט האמיתי. גם כאשר בגרסאות חדשות של
מודלים מסוימים קיימות גישות \textenglish{anchor-free}, העיקרון עדיין
חשוב להבנת הדרך שבה מערכות זיהוי מטפלות במיקום ובגודל.

לאחר החיזוי, מודל זיהוי עשוי להחזיר כמה תיבות שונות עבור אותו אובייקט.
לכן נהוג להשתמש ב־\textenglish{Non-Max Suppression (NMS)}, אלגוריתם
שבוחר בדרך כלל את התיבה בעלת ציון הביטחון הגבוה ביותר ומסיר תיבות
נוספות שחופפות לה מעל סף מוגדר.

בפרויקט שלי בחרתי להשתמש ב־\textenglish{YOLO26}, שהוא מודל
\textenglish{end-to-end}. לכן, בניגוד למודלים קלאסיים יותר, הוא יכול
להחזיר זיהוי סופי לכל אובייקט בלי להזדקק לשלב \textenglish{NMS} נפרד.
היתרון המעשי הוא תהליך חישוב פשוט ומהיר יותר, דבר שחשוב במיוחד כאשר
רוצים להריץ את המערכת גם על חומרה חלשה יחסית או על \textenglish{CPU} ולשמור על \textenglish{FPS} גבוה.

\subsubsection{מדדי הערכה: \textenglish{IoU}, \textenglish{Precision},
\textenglish{Recall} ו־\textenglish{mAP}}

כדי להעריך את איכות מודל הזיהוי משתמשים בכמה מדדים מרכזיים:
\begin{itemize}
\item \textbf{\textenglish{IoU (Intersection over Union)}:} יחס בין
שטח החפיפה של התיבה החזויה עם תיבת האמת, לבין שטח האיחוד שלהן. ערך
גבוה יותר מעיד על מיקום מדויק יותר.
\item \textbf{\textenglish{Precision}:} איזה חלק מן הזיהויים שהמודל
החזיר אכן נכונים.
\item \textbf{\textenglish{Recall}:} איזה חלק מן האובייקטים האמיתיים
אכן זוהו על ידי המודל.
\item \textbf{\textenglish{mAP}:} מדד מסכם, המשלב בין דיוק לשליפה
לאורך רמות ביטחון שונות, ונחשב לאחד המדדים המרכזיים בזיהוי אובייקטים.
\end{itemize}

\subsubsection{תיוג נתונים בפורמט \textenglish{YOLO}}

כדי לאמן מודל זיהוי אובייקטים, יש צורך בנתונים מתויגים. בכל תמונה יש
לסמן את מיקום האובייקט באמצעות \textenglish{Bounding Box} ולשייך לו
מחלקה. בפורמט \textenglish{YOLO} נהוג לייצג כל אובייקט באמצעות מזהה
מחלקה וארבעה ערכים מנורמלים: מרכז התיבה על ציר $x$, מרכז התיבה על ציר
$y$, רוחב התיבה וגובהה. בפרויקט זה השתמשתי ב־\textenglish{Roboflow}
לצורך תיוג הנתונים וניהולם.

\subsection{מעקב אובייקטים (\textenglish{Object Tracking})}

לאחר שהמערכת מזהה את הרובוטים בכל פריים בנפרד, יש צורך לקשר בין
הזיהויים לאורך זמן. זהו תפקידו של מעקב אובייקטים: לשמור על זהות עקבית
של אותו אובייקט גם כאשר הוא נע, נבלע מאחורי אובייקט אחר, או משנה את
מיקומו בתמונה.

\subsubsection{ההבדל בין \textenglish{Detection}
ל־\textenglish{Tracking}}

לעומת \textenglish{Detection} שמתייחס לכל פריים בפני עצמו, ולכן "שוכח" את מה
שקרה בפריים הקודם, \textenglish{Tracking} מחבר בין זיהויים
עוקבים ומקצה להם \textenglish{Track IDs} עקביים. בלי שלב המעקב, אפשר
היה לדעת היכן נמצאים הרובוטים בכל רגע, אך אי אפשר היה לדעת איזה רובוט
הוא אותו רובוט שראינו בפריים הקודם.

\subsubsection{מעקב מרובה־אובייקטים ו־\textenglish{Data Association}}

בבעיית \textenglish{Multi-Object Tracking} יש צורך לפתור בכל רגע בעיית
שיוך: איזה זיהוי חדש שייך לאיזו עקבה קיימת. תהליך זה נקרא
\textenglish{Data Association}. האלגוריתם מסתמך בדרך כלל על שילוב של
מיקום, מהירות משוערת, חפיפה בין תיבות ולעיתים גם דמיון חזותי.

בפרויקט שלי שלב זה חשוב במיוחד, משום שעל המגרש יש כמה רובוטים דומים,
ולעיתים הם מצטלבים, מסתירים זה את זה, או נעים במהירות גבוהה.

\subsubsection{\textenglish{ByteTrack} לעומת \textenglish{DeepSORT}}

שתי גישות נפוצות למעקב מרובה־אובייקטים הן:
\begin{itemize}
\item \textbf{\textenglish{DeepSORT}:} משתמש לא רק במיקום ובתנועה, אלא
  גם ב־\textenglish{Re-Identification} — כלומר בחתימה חזותית של
  האובייקט.  הדבר מאפשר התמודדות טובה יותר עם חסימות ארוכות, אך דורש
  חישוב נוסף ומודל חזותי נוסף.
\item \textbf{\textenglish{ByteTrack}:} מתמקד בעיקר בשיוך מבוסס תנועה,
  מיקום וחפיפה בין תיבות, אך באופן חכם מנצל גם זיהויים בעלי ביטחון נמוך.
  תכונה זו מועילה במיוחד במקרים של טשטוש תנועה, חסימה חלקית או זיהוי לא
  מושלם.
\end{itemize}

בפרויקט זה בחרתי ב־\textenglish{ByteTrack}, משום שהוא מספק איזון טוב
בין דיוק לבין מהירות, ומתאים ליישום שבו נדרש עיבוד וידאו יעיל
יחסית. בנוסף, אני חוסך אימון של מודל נוסף, וזיהוי הספרות כבר מבצע
בשבילי את ההתאמה מחדש ולא מתבלבל בין רובוטים דומים וויזואלית.

\subsubsection{\textenglish{ID Persistence}, חסימות
  ו־\textenglish{Re-Identification}}

אחד האתגרים המרכזיים במעקב הוא \textenglish{ID Persistence}: היכולת
לשמור על אותה זהות גם כאשר האובייקט מוסתר זמנית. כאשר רובוט חולף
מאחורי רובוט אחר, האלגוריתם צריך להחליט אם מדובר באותו אובייקט שחזר
להופיע או באובייקט חדש.

מערכות מסוימות פותרות זאת בעזרת \textenglish{Re-Identification}, כלומר
השוואת מאפיינים חזותיים של האובייקט בין הופעות שונות. בפרויקט שלי לא
נעשה שימוש ב־\textenglish{Re-Identification} מלא, משום שהבחירה הייתה
להעדיף פתרון מהיר ופשוט יותר חישובית. במקום זאת, המעקב נשען בעיקר על
התנועה, תוצאות הזיהוי והספרות שעל הבאמפר.

מבחינה תאורטית, מערכות \textenglish{Re-Identification} רבות נשענות על
למידת ייצוגים (\textenglish{Embeddings}) באמצעות \textenglish{Contrastive
Learning} או \textenglish{Triplet Loss}, כך שאובייקטים זהים יהיו קרובים
במרחב התכונות ואובייקטים שונים יהיו רחוקים זה מזה. בפרויקט שלי לא היה
צורך להוסיף רכיב כזה, משום שזיהוי הספרות מספק רמז זהות חזק יותר.

\subsubsection{דרישות \textenglish{FPS} וביצועים בזמן אמת}

במערכת \textenglish{Scouting} המיועדת לנתח וידאו של משחקים, המהירות
היא פרמטר קריטי. אם קצב העיבוד נמוך מדי, המערכת תתקשה להיות שימושית
בזמן אמת או אפילו בעיבוד מהיר לאחר המשחק. לכן יש יתרון למודלים ואלגוריתמים
קלים יחסית, שמסוגלים לעבוד בקצב \textenglish{FPS} גבוה, גם אם הם
מוותרים על חלק מן המורכבות של מודלים כבדים יותר.

זוהי אחת הסיבות המרכזיות לבחירה ב־\textenglish{YOLO}
וב־\textenglish{ByteTrack}: שני הרכיבים האלו תוכננו כך שיוכלו לספק
איזון טוב בין ביצועים לבין מהירות.

\subsubsection{כיוונים עתידיים: \textenglish{AprilTags} וכיול גאומטרי}

המערכת הנוכחית פועלת במישור התמונה הדו־ממדי, כלומר היא מנתחת את תנועת
הרובוטים יחסית לפיקסלים בתמונה. בעתיד ניתן לשלב
\textenglish{AprilTags} — סמנים חזותיים ידועים מראש — כדי לכייל את
המצלמה ולבצע מעבר מקואורדינטות של תמונה לקואורדינטות של המגרש האמיתי.

תוספת כזו עשויה לאפשר מדידה מדויקת יותר של מיקום, מהירות ומסלולים על
גבי המגרש עצמו, ולא רק בתוך מסגרת התמונה. עם זאת, היא דורשת תנאי צילום
מתאימים וכיול מצלמה, ולכן נחשבת לכיוון עתידי ולא לרכיב הכרחי במערכת
הנוכחית.
